% GENERATED FILE. Edit canonical lessons, then run npm run generate:books.
% canonical-source-hash: fnv1a64:fe946617f806cd29
% canonical-lessons: BN-C84-pulish, BN-C84-saikel, BN-C84-bajar, BN-C84-gram, BN-C84-shohor

\chapter{Police, Bicycle, Market, Village, City}
\label{ch:bn-police-bicycle-market-village-city}
\hlchaptermodality{84}

\begin{chapteropening}
\textbf{By the end of this chapter:} \emph{I can say the police, a bicycle, a market, a village and a city.}

Complete the last lesson of chapter 84: I can say the police, a bicycle, a market, a village and a city.
\end{chapteropening}

\section[\texorpdfstring{\bn{পুলিশ} (pulish)}{পুলিশ (pulish)}]{\bn{পুলিশ} (pulish) — the police}

\label{lesson:BN-C84-pulish}



\begin{quote}
Before the new one: say the Bengali for a temple, then the Bengali for a hotel.
\end{quote}

\subsection*{You'll want to know: \bn{পুলিশ}}
\textbf{\bn{পুলিশ}} — \emph{pulish} — \textquotedblleft{}the police\textquotedblright{}.

\begin{cousinweb}
\textbf{\bn{পুলিশ}} (\emph{pulish}) is a loan from English \emph{police}.
\end{cousinweb}

\begin{grammarlens}[title={What you've built}]
One more word for this chapter.
\end{grammarlens}

\subsection*{Your turn}
\begin{itemize}
\raggedright
  \item \emph{Say it:} \emph{pulish}
  \item \emph{Say it:} \emph{pulish}, once more
  \item \emph{Say it:} say \emph{pulish}
  \item \emph{Recall:} say the Bengali for a temple, then the Bengali for a hotel, then say \emph{pulish} again
\end{itemize}

\begin{tcolorbox}[breakable,colback=teal!4,colframe=teal!35!black,title={Before you move on}]
What does \bn{পুলিশ} mean? (\textquotedblleft{}The police\textquotedblright{}.) Say it once more.
\end{tcolorbox}
\section[\texorpdfstring{\bn{সাইকেল} (sāikel)}{সাইকেল (sāikel)}]{\bn{সাইকেল} (sāikel) — a bicycle}

\label{lesson:BN-C84-saikel}



\begin{quote}
Before the new one: say the Bengali for a hotel, then the Bengali for the police.
\end{quote}

\subsection*{You'll want to know: \bn{সাইকেল}}
\textbf{\bn{সাইকেল}} — \emph{sāikel} — \textquotedblleft{}a bicycle\textquotedblright{}.

\begin{cousinweb}
\textbf{\bn{সাইকেল}} (\emph{sāikel}) is a loan from English \emph{cycle}.
\end{cousinweb}

\begin{grammarlens}[title={What you've built}]
One more word for this chapter.
\end{grammarlens}

\subsection*{Your turn}
\begin{itemize}
\raggedright
  \item \emph{Say it:} \emph{sāikel}
  \item \emph{Say it:} \emph{sāikel}, once more
  \item \emph{Say it:} say \emph{sāikel}
  \item \emph{Recall:} say the Bengali for a hotel, then the Bengali for the police, then say \emph{sāikel} again
\end{itemize}

\begin{tcolorbox}[breakable,colback=teal!4,colframe=teal!35!black,title={Before you move on}]
What does \bn{সাইকেল} mean? (\textquotedblleft{}A bicycle\textquotedblright{}.) Say it once more.
\end{tcolorbox}
\section[\texorpdfstring{\bn{বাজার} (bājār)}{বাজার (bājār)}]{\bn{বাজার} (bājār) — a market}

\label{lesson:BN-C84-bajar}



\begin{quote}
Before the new one: say the Bengali for the police, then the Bengali for a bicycle.
\end{quote}

\subsection*{You'll want to know: \bn{বাজার}}
\textbf{\bn{বাজার}} — \emph{bājār} — \textquotedblleft{}a market\textquotedblright{}.

\begin{cousinweb}
\textbf{\bn{বাজার}} (\emph{bājār}) is a loan, usually traced to Persian \emph{bāzār}, \textquotedblleft{}market\textquotedblright{}.
\end{cousinweb}

\begin{grammarlens}[title={What you've built}]
One more word for this chapter.
\end{grammarlens}

\subsection*{Your turn}
\begin{itemize}
\raggedright
  \item \emph{Say it:} \emph{bājār}
  \item \emph{Say it:} \emph{bājār}, once more
  \item \emph{Say it:} say \emph{bājār}
  \item \emph{Recall:} say the Bengali for the police, then the Bengali for a bicycle, then say \emph{bājār} again
\end{itemize}

\begin{tcolorbox}[breakable,colback=teal!4,colframe=teal!35!black,title={Before you move on}]
What does \bn{বাজার} mean? (\textquotedblleft{}A market\textquotedblright{}.) Say it once more.
\end{tcolorbox}
\section[\texorpdfstring{\bn{গ্রাম} (grām)}{গ্রাম (grām)}]{\bn{গ্রাম} (grām) — a village}

\label{lesson:BN-C84-gram}



\begin{quote}
Before the new one: say the Bengali for a bicycle, then the Bengali for a market.
\end{quote}

\subsection*{You'll want to know: \bn{গ্রাম}}
\textbf{\bn{গ্রাম}} — \emph{grām} — \textquotedblleft{}a village\textquotedblright{}.

\begin{cousinweb}
\textbf{\bn{গ্রাম}} (\emph{grām}) is Sanskrit \emph{grāma}, \textquotedblleft{}troop, village\textquotedblright{}, taken over whole as a learned loan, a tatsama (CDIAL 4368). The word inherited from it is \emph{gāṁ}, \textquotedblleft{}village\textquotedblright{}, which CDIAL lists.
\end{cousinweb}

\begin{grammarlens}[title={What you've built}]
One more word for this chapter.
\end{grammarlens}

\subsection*{Your turn}
\begin{itemize}
\raggedright
  \item \emph{Say it:} \emph{grām}
  \item \emph{Say it:} \emph{grām}, once more
  \item \emph{Say it:} say \emph{grām}
  \item \emph{Recall:} say the Bengali for a bicycle, then the Bengali for a market, then say \emph{grām} again
\end{itemize}

\begin{tcolorbox}[breakable,colback=teal!4,colframe=teal!35!black,title={Before you move on}]
What does \bn{গ্রাম} mean? (\textquotedblleft{}A village\textquotedblright{}.) Say it once more.
\end{tcolorbox}
\section[\texorpdfstring{\bn{শহর} (shôhôr)}{শহর (shôhôr)}]{\bn{শহর} (shôhôr) — a city}

\label{lesson:BN-C84-shohor}



\begin{quote}
Before the new one: say the Bengali for a market, then the Bengali for a village.
\end{quote}

\subsection*{You'll want to know: \bn{শহর}}
\textbf{\bn{শহর}} — \emph{shôhôr} — \textquotedblleft{}a city\textquotedblright{}.

\begin{cousinweb}
\textbf{\bn{শহর}} (\emph{shôhôr}) is a loan, usually traced to Persian \emph{shahr}, \textquotedblleft{}city\textquotedblright{}.
\end{cousinweb}

\begin{grammarlens}[title={What you've built}]
One more word for this chapter.
\end{grammarlens}

\subsection*{Your turn}
\begin{itemize}
\raggedright
  \item \emph{Say it:} \emph{shôhôr}
  \item \emph{Say it:} \emph{shôhôr}, once more
  \item \emph{Say it:} say \emph{shôhôr}
  \item \emph{Recall:} say the Bengali for a market, then the Bengali for a village, then say \emph{shôhôr} again
\end{itemize}

\begin{tcolorbox}[breakable,colback=teal!4,colframe=teal!35!black,title={Before you move on}]
What does \bn{শহর} mean? (\textquotedblleft{}A city\textquotedblright{}.) Say it once more.
\end{tcolorbox}
